\section{Computing Chow Groups}

\blemm

    If $X$ is a variety of dimension $n$ then $\Ch_n(X)=\bZ[X]$.

\elemm

\bproof

    $X$ is its only subvariety of dimension $n$, and it has no $(n+1)$-dimensional subvarieties and
    so $\rat_n(X)=0$.
    \qed

\eproof

\subsection{Affine Space}

\bprop[name={chowaffine}]

    The Chow group of $\bA^n$ is the free Abelian group generated by $[\bA^n]$.
    Meaning
    $$ \Ch(\bA^n) = \bZ[\bA^n] . $$

\eprop

\bproof

    Let $Y\subseteq\bA^n$ be a proper subvariety, and let $z_1,\dots,z_n$ be coordinates on $\bA^n$
    such that the origin doesn't lie in $Y$.
    Define
    $$ W^\circ = \set{(t,tz)\in(\bA^1-0)\times\bA^n}[z\in Y] = \vof\set{f(z/t)}[f(z)\in\vof(Y)] , $$
    so that the fiber of $W^\circ$ over any $t\in\bA^1-0$ is $tY$.
    And $W^\circ$ is irreducible as it is the image of $(\bA^1-0)\times Y$.
    Define $W\subseteq\bP^1\times\bA^n$ be the closure of $W^\circ$ in $\bP^1\times\bA^n$, which is irreducible
    as $W^\circ$ is.

    Because $0\notin Y$, there is a polynomial $g(z)\in\vof(Y)$ with nonzero constant term $c$.
    Define $G(t,z)=g(z/t)$ on $(\bA^1-0)\times\bA^n$, and this extends to a automorphism of
    $(\bP^1-0)\times\bA^n$ whose value on $\infty\times\bA^n$ is the constant $c$.

    Thus the fiber of $W$ over $t=1$ is $Y$, and the fiber of $W$ over $t=\infty$ is empty.
    This is because $G$ vanishes on $W$, and when $t=\infty$ it does not vanish.
    Thus $[Y]=[W^{-1}0]-[W^{-1}\infty]$ is then zero.
    \qed

\eproof

\subsection{Mayer-Vietoris}

If $Y\subseteq X$ is a closed subscheme, then there is an obvious inclusion $Z_k(Y)\inj Z_k(X)$.
Furthermore this maps $\div(r)$ for $r\in K(W)^\times$ for $W\subseteq Y$ to $\div(r)\in\rat_k(X)$.
Thus this descends into a map $\Ch_k(Y)\to\Ch_k(X)$.

Conversely if $U\subseteq X$ is open, then there is an obvious map $Z_k(X)\to Z_k(U)$ which intersects a
variety with $U$, and this preserves rational equivalence as well.
And so it descends into a map $\Ch_k(X)\to\Ch_k(U)$.

\bthrm[name={excision}, title={Excision}]

    If $Y\subseteq X$ is a closed subscheme and $U=X-Y$ its complement, then the following sequence is
    exact
    $$ \Ch(Y) \to \Ch(X) \to \Ch(U) \to 0 . $$

\ethrm

\bproof

    Denote the inclusions by $i\colon Y\to X,j\colon U\to X$, and the induced maps by $i_*,j^*$.
    Because a subvariety $V$ of $U$ extends to a subvariety $\cl V\subseteq X$, the following sequence is
    exact
    $$ Z_kY \xtos{i_*} Z_kX \xtos{j^*} Z_kU \to 0 . $$
    Thus we need only show exactness at $\Ch_k(X)$.
    Suppose that $\alpha\in Z_k(X)$'s pullback $j^*\alpha$ is rationally equivalent to zero.
    Thus
    $$ j^*\alpha = \sum[\div r_i] $$
    for $r_i\in K(W_i)^\times$ where $W_i\subseteq U$ are subvarieties.
    Because $K(W_i)=K(\cl W_i)$, $r_i$ corresponds to a rational function $\cl r_i\in K(\cl W_i)$.
    Furthermore we see that $j^*[\div(\cl r_i)]=\div(r_i)$ so that
    $$ j^*\parens{\alpha-\sum[\div\cl r_i]} = 0 \in Z_k(U) $$
    and thus $\alpha-\sum[\div\cl r_i]=i_*\beta$ for $\beta\in Z_kY$, showing exactness.
    \qed

\eproof

\bthrm[name={mayerviet}, title={Mayer-Vietoris}]

    Let $X_1,X_2\subseteq X$ be closed subschemes, and let $i_1,i_2$ denote the inclusions of $X_j$ into $X$.
    Then there is a right exact sequence
    $$ \Ch(X_1\cap X_2) \xtos\phi \Ch(X_1)\oplus\Ch(X_2) \xtos\psi \Ch(X_1\cup X_2) \to 0 . $$
    Here $\phi=((i_1)_*,-(i_2)_*)$ and $\psi=((i_1)_*,(i_2)_*)$.

\ethrm

\bproof

    We may assume $X=X_1\cup X_2$.
    Then we have an exact sequence
    $$ Z_k(X_1\cap X_2) \xtos\phi Z_k(X_1)\oplus Z_k(X_2) \xtos\psi Z_k(X) \to 0 . $$
    And if $\psi([A],[B])=[A]+[B]=\sum[\div(r_i)]$ then $\psi\parens{[A]-\sum[\div r_i],[B]}=0$ and thus
    there is a $\beta\in Z_k(X_1\cap X_2)$ such that $\phi(\beta)=\parens{[A]-\sum[\div r_i],[B]}$ and
    thus they are rationally equivalent.
    \qed
    \TODO{Check}

\eproof

\bcoro[name={chowquasiaffine}]

    If $U\subseteq\bA^n$ is nonempty then $\Ch(U)=\Ch_n(U)=\bZ[U]$.

\ecoro

\bproof

    If we let $Y=\bA^n-U$ then the map $\Ch(Y)\to\Ch(\bA^n)$ is the zero map, and thus by
    \refmath{excision} $\Ch(\bA^n)\to\Ch(U)$ is a bijection as desired.
    \qed

\eproof

\subsection{Affine Stratification}

\bdefn

    A {\emphcolor stratification} of a topological space $X$ is a finite collection of disjoint irreducible locally-closed subspaces $U_i$ of $X$
    which cover $X$ and such that if $\cl U_i\cap U_j$ is nonempty then $U_j$ is contained in $\cl U_i$.
    The $U_i$ are called the {\emphcolor strata} and their closures $Y_i=\cl U_i$ are the {\emphcolor closed strata}.

\edefn

We can obtain the strata from the closed strata:
$$ U_i = Y_i - \bigcup_{Y_j\subsetneq Y_i}Y_j . $$

\bnote

    Note that $Y_i=Y_j$ iff $i=j$.
    Indeed, we know that $Y_i$ are irreducible as $U_i$ are, and so if $U_i,U_j\subseteq Y_i=Y_j$ then because they are open
    we have that they intersect non-trivially, so $i=j$.

    Now if $U_i\cap Y_j\neq\emptyset$ then $U_i\subseteq Y_j$ and thus $Y_i\subseteq Y_j$ and in particular $Y_j$ is not a proper
    subset of $Y_i$.
    Thus $U_i$ does not intersect any of the $Y_j\subsetneq Y_i$, and so $U_i\subseteq Y_i-\bigcup_{Y_j\subsetneq Y_i}Y_j$.
    Conversely if $x\in Y_i-U_i$ then $x\in U_j$ for some $j\neq i$ and thus $U_j\cap Y_i\neq\emptyset$ and thus $Y_j\subseteq Y_i$.
    Because $i\neq j$ we have $Y_j\subsetneq Y_i$, so $x\in\bigcup_{Y_j\subsetneq Y_i}Y_j$.

\enote

We form a partial order on the strata of a stratification by having $i\leq j$ iff $U_i\subseteq Y_j$.

For instance if we have a flag of closed irreducible subsets $Y_0\subsetneq Y_1\subsetneq\cdots\subsetneq Y_n=X$, then
they form the closed strata of a stratification.

\bdefn

    A stratification of a scheme $X$ with strata is $U_i$ is {\emphcolor (quasi-)affine} if each $U_i$ is (quasi-)affine.

\edefn

For instance $Y_i=\bP^i\subseteq\bP^n$ forms an affine stratification of $\bP^n$, as $U_i=\bP^i-\bP^{i-1}\cong\bA^i$.

\bprop[name={stratagenerate}]

    If a scheme $X$ has a quasi-affine stratification, then $\Ch(X)$ is generated by the classes of the closed strata
    as an Abelian group.

\eprop

\bproof

    The proof is by induction on the number of strata $U_i$.
    If $n=1$ then $X$ is quasi-affine, and the result follows from \refmath[corollary]{chowquasiaffine}.

    Now let $U_0$ be a minimal strata (relative to the partial order on strata).
    Because $Y_0$ is the union of strata, and $U_0$ is minimal, we must have that $U_0$ is already closed.
    Thus $U=X-U_0$ is stratified by the other $U_i$, and by induction $\Ch(U)$ is generated by the classes of these $U_i$.
    By \refmath{excision} there is an exact sequence
    $$ \bZ[U_0] = \Ch(U_0) \to \Ch(X) \to \Ch(U) \to 0 . $$
    Because the images of $[Y_i]\in\Ch(X)$ generate $\Ch(U)$, we see that $[Y_0],[Y_i]$ generate $\Ch(X)$ as desired.
    \qed

\eproof

\bnote

    The closed strata of a quasi-affine stratification may be zero in the Chow group.
    For instance take the quasi-affine stratification of $\bA^1$ consisting of a point and its complement.
    By \refmath[proposition]{chowaffine} we know that the class of a point in an affine Chow group is zero.

    But when the stratification is affine, a theorem of Totaro [2014] says that the closed strata actually form a basis
    of $\Ch(X)$.

\enote

\subsection{Projective Space}

Consider the map $\deg\colon Z_0(\bP^n)\to\bZ$ defined by mapping a closed point $p\in\bP^n$ to $[\kappa(p):k]$, where $\kappa(p)$ is the residue field
of $p$, and $[\kappa(p):k]$ is the degree of the field extension $\kappa(p)/k$.
If $k$ is algebraically closed this just maps $p$ to $1$.
We extend this to the rest of $Z(\bP^n)$ by mapping all cycles of positive pure dimension to zero.
It can be shown that the degree of a principal divisor $\div f$ is zero, and thus $\deg$ defines a map $\Ch(\bP^n)\to\bZ$.

\bnote

    This is a special case of a more general result.
    If $f\colon X\to Y$ is a proper morphism of varieties, then $f$ maps subvarieties of $X$ to subvarieties of $Y$.
    Furthermore for a variety $V\subseteq X$, there is a natural embedding $K(f(V))\subseteq K(V)$ which is finite when $\dim V=\dim f(V)$.
    Then define $\deg_f(V)=[K(V):K(f(V))]$ when the dimensions are equal and zero otherwise.

    Then we define
    $$ f_*\colon Z_kX \to Z_kY,\qquad f_*[V] = \deg_f(V)[fV] . $$
    It can be shown (e.g.\ Fulton) that this maps divisors to divisors, and thus descends to a morphism
    $$ f_*\colon\Ch_k(X) \to \Ch_k(Y) . $$

    In the case that $X$ is proper over $k$ (i.e.\ $X\to\spec k$ is a proper morphism), then the map $X\to\spec k$ defines a map
    $$ \deg\colon\Ch(X) \to \Ch(\spec k) \cong \bZ[\spec k] \cong \bZ , $$
    which maps a point $[p]$ to $[\kappa(p):k]\in\bZ$.

\enote

\blemm

    We have $\Ch_0(\bP^n)=\Ch^n(\bP^n)=\bZ$.

\elemm

\bproof

    The closed points of $\bP^n$ are all rationally equivalent.
    For two closed points $a,b\in\bP^n$ we may parameterize $\bP^n$ so that the points lie on a copy of $\bP^1$ and are of the form
    $[1:0]$ and $[0:1]$.
    Then $r(x/y)=x/y$ is an invertible rational function whose divisor is $\set a-\set b$.

    So $\Ch_0(\bP^n)$ is generated by $[p]$, and the degree of $n[p]$ is $n\neq0$, so this is a basis.
    \qed

\eproof

\bthrm[name={chowprojective}]

    Let $H\subseteq\bP^n$ be a hyperplane and let $\zeta\in\Ch^1(\bP^n)$ be its class.
    Then as rings
    $$ \Ch(\bP^n) = \frac{\bZ[\zeta]}{(\zeta^{n+1})} . $$

\ethrm

\bproof

    Let $\bP^0\subsetneq\cdots\subsetneq\bP^n$ be an affine stratification, with strata $U_i=\bP^i-\bP^{i-1}$.
    Then by \refmath[proposition]{stratagenerate}, $\bP^{n-k}$'s class generates $\Ch^k(\bP^n)$.
    Thus any $(n-k)$-plane $L$ generates $\Ch^k(\bP^n)$.
    Let $M$ be a general $k$-plane, so that $L$ and $M$ intersect transversally in one point.
    Thus $n[L]\cdot[M]=n\in\Ch^n(\bP^n)=\bZ$ is nonzero, and thus $[L]$ freely generates $\Ch^k(\bP^n)$.
    That is, $\Ch^k(\bP^n)\cong\bZ$.

    Because $L$ is the transverse intersection of $k$ hyperplanes, $[L]=\zeta^k$.
    Thus the obvious map $\bZ[\zeta]\to\Ch(\bP^n)$ has kernel $(\zeta^{n+1})$ and is surjective, as desired.
    \qed

\eproof

\bcoro[name={projvarinchow}]

    If $X\subseteq\bP^n$ has codimension $k$ and degree $d$, then $[X]=d\zeta^k$.

\ecoro

\bproof

    We know that $[X]=m\zeta^k$ for some $m\in\bZ$ by the above theorem.
    $X$ intersects a general $k$-plane transversally in $d$ points, so that
    $$ d = \deg([X]\zeta^{n-k}) = \deg(m\zeta^n) = m $$
    as $\zeta^n$ is a single point.
    So $[X]=d\zeta^k$ as desired.
    \qed

\eproof

\bcoro[title={B\'ezout}]

    Let $X_1,\dots,X_k\subseteq\bP^n$ be subvarieties of codimension $c_1,\dots,c_k$ with $\sum c_i\leq n$ and the $X_i$ intersect
    generically transversally, then
    $$ \deg(X_1\cap\cdots\cap X_k) = \prod\deg(X_i) . $$

\ecoro

\bproof

    Define $d_i=\deg X_i$, so we have
    $$ \deg(X_1\cap\cdots\cap X_k) = \deg([X_1]\cdots[X_k]\zeta^{n-c}) =
    \deg(d_1\zeta^{c_1}\cdots d_n\zeta^{c_n}\zeta^{n-c}) = \deg(d_1\cdots d_n\zeta^n) = d_1\cdots d_n $$
    as desired.
    \qed

\eproof

